\chapter{Conclusioni e sviluppi futuri}

In questo lavoro di tesi è stato affrontato il problema della classificazione
di file Windows PE come benigni o malevoli attraverso la loro
rappresentazione sotto forma di immagini e tramite l'utilizzo di una rete
ResNet-50. In particolare, sono state confrontate le strategie di
mapping LINEAR, ZIGZAG e SERPENTINE, considerando sia immagini RGB
sia immagini in scala di grigi e le due diverse strategie relative alle
dimensioni, a dimensione fissa e ad altezza variabile.
Nel complesso, i risultati ottenuti hanno mostrato prestazioni
elevate per tutte le configurazioni considerate, evidenziando
risultati migliori per le configurazioni ad altezza variabile
rispetto alle corrispondenti configurazioni a dimensione fissa.
In particolare, la migliore overall accuracy è stata ottenuta
dalla configurazione SERPENTINE RGB ad altezza variabile, con un
valore pari al 97,61\%.
L'approccio basato sulla rappresentazione dei file PE come immagini
e sulla loro classificazione mediante ResNet-50 ha permesso di
ottenere buone prestazioni sul dataset considerato, mostrando inoltre
come la scelta delle dimensioni delle immagini possa avere un ruolo
importante nelle prestazioni del modello.

Come possibile sviluppo futuro, l'analisi potrebbe essere estesa
utilizzando dataset differenti, in modo da verificare se le
prestazioni rimangano coerenti anche su altri insiemi di file PE.
Un ulteriore sviluppo potrebbe riguardare il confronto
\textcolor{red}{della accuratezza della classificazione ottenuta con la} rete
ResNet-50 \textcolor{red}{con l'accuratezza che si potrebbe ottenere usando} altre architetture della stessa
famiglia, caratterizzate da una diversa profondità, in modo da
valutare l'influenza del numero di livelli sulle prestazioni ottenute.
Potrebbero inoltre essere valutate diverse modalità per determinare
le dimensioni delle immagini, ad esempio calcolandole in funzione
della lunghezza in byte del singolo file PE, invece di utilizzare
una larghezza comune determinata a partire dalla lunghezza media in
byte dei file del dataset.
\textcolor{red}{Nella versione attuale della tesi, la fase di pre-elaborazione}
delle immagini, \textcolor{red}{compresa la fase di ridimensionamento (resize)}
\textcolor{red}{delle immagini, sono state, entrambe, eseguite} prima di fornire
le immagini in input alla ResNet-50, mentre potrebbe essere interessante valutare come cambiano
le prestazioni utilizzando l'AutoImageProcessor di Hugging Face
\cite{huggingface_autoimageprocessor} per la preparazione delle
immagini.
Questi possibili sviluppi permetterebbero di proseguire l'analisi e
approfondire \textcolor{red}{, in futuro, come la} rappresentazione dei file PE come
immagini \textcolor{red}{possa contribuire a migliorare l'accuratezza della} classificazione di file benigni e malevoli.